\chapter{绪论}\label{chap:introduction}

\section{研究背景与意义}

\subsection{研究背景}

近年来，国产处理器与基础软件生态建设持续推进，LoongArch 作为面向自主可控需求设计的
指令集架构，其基础架构、向量扩展和 ABI 规范已形成较完整的公开文档体系\citep{LoongArchISA2023,
LoongArchISAv2,LoongArchELFABI2025}。与此同时，移动操作系统与应用生态已成为处理器
软件支持能力的重要体现。Android 长期占据全球移动操作系统主导地位\citep{StatCounter2025}，
其运行时系统直接影响应用启动速度、交互时延和整体执行效率。因此，在国产处理器平台上完善
Android 核心运行环境，关键不只在于把代码移植过来，更在于让应用生态真正可用，并把性能控制在可接受范围内。

Android Runtime（ART）是 Android 应用执行的核心运行时。ART 承担 DEX 字节码解释执行、
即时编译、安装期预编译、JNI 调用桥接、类加载和运行时支持等任务\citep{AndroidRuntime2024,
ArtJitCompiler2024,AospArtAndroid15R1}。相较于早期 Dalvik，ART 采用解释执行、Just-In-Time（JIT）
即时编译与 Ahead-Of-Time（AOT）预编译相结合的混合执行模式，在启动开销、生成代码质量与运行时
资源占用之间寻求平衡\citep{ArtJitCompiler2024,Paleczny2001,Suganuma2001}。这一机制
决定了 ART 的平台适配不是补一个后端即可，而是要同时处理解释器、编译器后端、JNI 编译器和运行时
入口点等多个子系统。

然而，截至本文写作时，AOSP 官方主线的 ART 源码目录中仍未见面向 LoongArch 的执行引擎
后端实现\citep{AospArtAndroid15R1,AospArtMain2026}。对于新架构平台而言，
即使 Linux 内核、Bionic 和系统框架层已经具备基础可运行能力，若缺少 ART 层支持，Android
应用生态仍难以真正落地。特别是 ART 中与体系结构密切相关的解释器入口、Nterp 汇编解释器、
JIT 代码生成、JNI 调用约定和快速路径入口点等模块，均需要针对目标架构进行专门适配。此外，
在完成“能运行”的功能目标之后，仍需进一步解决“能正确”和“能更快”的问题，即通过系统回归
验证保证语义正确性，并通过性能分析与优化降低新平台上的执行开销。

\subsection{研究意义}

围绕 LoongArch 平台 ART 执行引擎的适配与优化开展研究，意义主要体现在以下三个方面。

第一，该研究有助于补齐国产处理器上的 Android 运行时关键能力。相比通用用户态程序移植，
运行时系统直接决定 Java 与 Kotlin 应用能否稳定、高效地在目标平台执行，因此 ART 的移植结果
直接影响 LoongArch 平台 Android 应用生态能否落地。

第二，该研究能够积累新架构平台运行时系统适配的方法与经验。ART 并非单纯的解释器或编译器，
而是同时包含解释执行、JIT/AOT 编译、JNI 桥接、异常处理、调试信息和运行时快路径的复杂系统。
对其进行平台化适配，可以为新架构上的执行引擎移植积累可供借鉴的经验。

第三，该研究为 LoongArch 平台上的性能优化研究建立了实验基础。只有在执行引擎基础能力补齐并经
系统验证后，围绕具体工作负载的性能瓶颈分析与优化才具备可靠前提。本文在完成移植适配后进一步
开展基准测试驱动的优化研究，并在统一实验条件下评估不同优化方向的有效性与适用范围。

\section{国内外研究现状}

\subsection{Android运行时与混合编译研究}

Android 运行时经历了从 Dalvik 向 ART 的演进过程。随着应用规模扩大，仅依赖解释执行已难以
满足性能需求，运行时中逐步引入 JIT 与 AOT 混合编译机制\citep{ArtJitCompiler2024,Shi2008,
Wang2011AOTAndroid}。
已有研究显示，移动端运行时的关键在于在热点增长节奏与编译成本之间找到合适平衡，
而非单纯追求更激进的编译等级\citep{Paleczny2001,Suganuma2001,Kotzmann2008}。

围绕 Android/ART 平台 JIT 行为本身，也已有工作从混合编译协作、代码缓存复用和低开销运行时观测等角度
开展优化研究。Hybrid JIT 工作比较了不同 JIT 形态的协作方式，ShareJIT 关注跨进程代码缓存复用，
ARTist 说明了基于运行时编译链进行插桩与观测的可行性，Hubble 则展示了面向 Android 的低开销
方法级跟踪机制\citep{Perez2012HybridJIT,Xu2018ShareJIT,Backes2017ARTist,Luo2022Hubble}。
这些研究说明，ART 的性能优化通常需要
紧密结合工作负载生命周期和热点路径特征，不能简单依赖单一静态阈值或全局默认策略。

进一步看，近年来 Android/ART 研究的重点已经不只限于“是否采用混合编译”，而是逐步扩展到代码缓存共享、
运行时观测和编译链插桩等方向。ShareJIT 通过跨进程共享 JIT 代码缓存，试图降低重复编译与代码缓存占用
\citep{Xu2018ShareJIT}；Hubble 利用 Android 运行时环境，在方法级 tracing 的同时将可容忍开销控制在
面向生产环境的范围内\citep{Luo2022Hubble}；ARTist 则表明，基于 ART 编译链进行运行时插桩与安全分析
是可行的\citep{Backes2017ARTist}。这些工作共同说明，既有 Android 架构上的执行引擎研究，已经从
“如何生成代码”进一步扩展到“如何复用代码、如何观察代码、如何分析代码”。

\subsubsection{2023-2025年的最新进展：编译器优化与垃圾回收}

2023--2025 年间，ART 官方公开材料显示，其优化重点已从单一编译策略选择，进一步扩展到编译器自身效率
与代码体积控制。2025 年的 Android Developers Blog 文章指出，ART 团队在不牺牲峰值内存和生成代码质量
的前提下，将编译时间缩短了 18\%\citep{GoogleART2025}。2023 年的另一篇官方文章则总结了代码体积优化结果，
其中写屏障消除（Write Barrier Elimination）、返回指令合并等改进合计带来了约 9.3\% 的编译代码体积缩减
\citep{GoogleART2023}。

垃圾回收（GC）方面，AOSP 文档表明 Concurrent Copying 在 Android 10 及以上版本默认以分代模式运行，
通过更频繁地回收年轻对象来提升吞吐并推迟 full-heap GC\citep{GoogleART2024}。与此同时，ISMM 2024 的研究
指出，ART 上垃圾回收开销在不同工作负载之间差异显著，在其评测中下界可从 2\% 到 51\% 不等，因此移动运行时
上的 GC 评估需要配套的方法学约束与基准设计\citep{GCEvaluation2024}。

这些进展表明，ART 的优化已经同时覆盖编译速度、代码体积与内存管理行为；相比早期更关注“是否编译、何时编译”，
近年的问题意识更强调“如何更高效地编译，以及如何更可控地管理运行时开销”。

本文与上述工作的关注点并不相同。ShareJIT、Hubble 与 ARTist 都建立在既有架构平台执行能力已经成立的
前提上，研究重点分别是缓存共享、低开销观测与运行时插桩；而本文首先要回答的是，在 LoongArch 这一
新架构平台上，ART 的解释器、Nterp、JIT、JNI 和运行时入口点能否被系统性补齐，并形成可运行、可验证、
可比较的执行基础。只有这一前提成立，后续针对 LoongArch 平台的性能分析与优化结论才具备研究意义。

\subsection{解释器与执行引擎优化研究}

在执行引擎层面，解释器性能长期是虚拟机研究中的重要问题。经典研究指出，解释器的分派结构与
分支预测行为会显著影响整体吞吐\citep{Ertl2003}；与此同时，JIT 编译虽能提升代码质量，但也会
引入编译线程开销与调试信息等固定成本\citep{Arnold2005,Fink2003,Lameed2013}。因此，执行引擎
优化的关键在于平衡解释器驻留时间与 JIT 编译触发时机，而非单纯追求更早或更激进的编译行为。
关于解释器分派机制和 JIT 分层策略的详细技术分析，将在第 2 章中进一步展开。

对于字符串与运行时热路径优化，除解释器与编译器层面的优化外，运行时热路径中的具体操作同样值得关注。
已有研究显示，Java 字符串对象的存储布局、复制方式和字符编码
转换路径会对内存占用与执行效率产生持续影响\citep{Haubl2008OptimizedStrings,
Haubl2010CompactStrings}。特别是在大量短字符串构造、拼接与替换场景中，
热路径中的逐元素搬运和额外控制开销容易成为稳定瓶颈。因此，围绕字符串处理路径开展局部、低开销
的微优化，仍然是运行时系统中具有现实价值的方向。

\subsubsection{SIMD向量化与相关工具链现状}

在 SIMD 向量化方面，Arm 官方资料将 Neon 视为固定宽度的 SIMD 扩展，将 SVE/SVE2 视为面向 AArch64 的
可伸缩向量扩展，并强调 SVE2 既可以由编译器从 C/C++ 循环自动生成，也可以通过 intrinsics 或汇编显式使用
\citep{ARMSIMD2023,ARMCompiler2024}。这说明近年的关注点不仅是“处理器是否提供 SIMD 指令”，还包括
编译器自动向量化能力以及 intrinsics 级别的可用性。

对 LoongArch 而言，LLVM 官方发布说明显示，17.0.1 版本的 LoongArch 后端已经加入 LSX、LASX 等 ISA
扩展的汇编/反汇编支持，并持续改进代码生成\citep{LoongArchLLVM2023}。与此同时，社区已出现面向
LoongArch 的运行时代码生成工具，如 xbyak\_loongarch64 项目\citep{XbyakLoongArch2024}；V8 文档
也将 Loong64 列为独立的非官方支持移植端口\citep{V8LoongArch2024}。这些材料共同说明，LoongArch 的
SIMD 与 JIT 相关基础设施已经形成雏形，但整体生态仍处于持续补齐和演进阶段。

\subsection{LoongArch平台JVM与ART移植研究}

在国产处理器平台上的虚拟机与 Android 运行时移植方面，已有研究已从早期系统整体移植推进到围绕
LoongArch 平台执行环境关键组件的适配探索。黄佳的工作完成了基于 LoongArch64 指令集的 Java
虚拟机移植与优化，验证了在新架构平台上补齐汇编器、解释器和编译器后端的可行性\citep{Huang2021}。
杜泽旭面向龙芯 2K 平台完成了 Android 系统移植与优化，涉及 ART 运行时与图形子系统适配
\citep{Du2020Loongson2KAndroid}；更早期面向龙芯平台的 Android 系统移植研究也说明，系统整体移植、
虚拟机适配与性能验证需要协同推进\citep{Lei2012AndroidLoongson}。这些工作为本文提供了直接参考，
也说明国产平台上的 Android 运行时适配并非完全空白。但对于面向 AOSP 15 的较新版本 ART，如何在固定
对照版本下组织工作负载驱动的优化研究，并用统一实验设计支撑正式性能结论，现有公开工作仍讨论较少。

与黄佳的工作相比，本文的研究对象已经从 OpenJDK/HotSpot 转向 AOSP 15 的 ART。前者重点讨论
HotSpot 的汇编器、模板解释器和 C2 编译器移植与优化\citep{Huang2021}；而本文面对的是同时包含
C++ 解释器、Nterp、分层 JIT、JNI 编译器和运行时入口点的 Android 执行引擎，因此问题重心从
单一 JVM 执行部件移植，扩展为多个执行子系统之间的协同适配与边界校正。

与杜泽旭的工作相比，本文也不再以 Android 系统整体移植为主线。杜泽旭面向龙芯 2K 平台和
Android 5.1.1，重点讨论系统整体移植、ART 的 MIPS64 后端支持以及图形子系统优化
\citep{Du2020Loongson2KAndroid}；本文则基于 LoongArch64 平台和 AOSP 15，把 ART 执行引擎本身
作为独立研究对象，更强调执行能力构建、统一实验设计以及平台相关优化问题的筛选与验证。

因此，本文与已有工作的差异并不只在于平台或版本更新，更在于研究目标已经从“让系统在国产平台上运行起来”
进一步收缩为“让执行引擎能力建立起来、让性能结论具备统一证据基础，并在此基础上识别平台相关的优化方向”。

因此，面向 AOSP 15 的 LoongArch 平台 ART 执行引擎适配与优化仍然具有明显研究空间。尤其是在
新版本 ART 已同时包含 C++ 解释器、Nterp、JIT/AOT 编译、JNI 编译和更复杂运行时支撑机制的背景下，
如何完成新版本源码适配，并在系统验证基础上开展更有针对性的优化研究，仍是一个值得回答的问题。

\section{本文研究内容与主要贡献}

\subsection{研究问题}

本文围绕一个核心问题和两个递进子问题展开研究。

核心问题：在 AOSP 15 这一较新版本 ART 上，LoongArch 平台能否建立完整、可运行的执行引擎能力链，
以及这一能力链应如何组织。该问题对应第 3 章，重点在于解释器、Nterp、JIT 编译器后端、JNI 编译器和
运行时快速路径之间的协同关系。

支撑问题：在适配完成后，如何建立既能证明系统正确性、又能支撑性能结论的统一实验验证体系。该问题
对应第 4 章，重点在于组织目标机回归、libcore 测试、基准测试评测以及统一的正式对比数据集。

验证问题：在统一实验条件和固定对照版本下，哪些平台相关优化方向真正值得进入正文，并能够形成强弱有别的
性能结论。该问题对应第 5 章和第 6 章：第 5 章负责问题建模、候选方向筛选与代表性案例分析，第 6 章给出
LoongArch64 最小可用 LSX 浮点 SIMD 路径这一全文主结论，并通过重复测量和边界比较约束结论强度。

\subsection{研究内容}

本文以 AOSP 15 的 ART 为研究对象，围绕 LoongArch 平台执行引擎的适配与优化开展工作，主要包括
以下四个方面。

第一，完成 LoongArch 平台 ART 执行能力构建。本文以执行引擎相关模块为重点，系统
补齐汇编器、C++ 解释器、Nterp 汇编解释器、JIT 编译器后端、JNI 编译器和运行时快速路径等核心
能力，使 ART 在 LoongArch 平台上具备完整执行基础。

第二，建立面向目标设备的统一实验设计与证据组织。本文围绕目标机回归、libcore 验证和基准测试评测
组织统一实验流程，并将正文使用的数据抽象为 \texttt{Snapshot-A} 至 \texttt{Snapshot-E} 五类
实验快照，以区分正式结果、统计验证、启动期案例和补充消融材料。

第三，围绕执行期与启动期路径组织代表性局部优化案例。本文将 JIT 进入时机与运行时热路径数据搬移
分别作为控制流与数据流方向的代表性瓶颈，并将 DEX 校验阶段链作为正交启动期案例，分析其机制、
收益边界与适用范围。

第四，补齐 LoongArch64 最小 LSX 浮点向量 lowering 链路，并以此形成全文主结论。本文围绕规则
浮点数值内核，给出从 Loop Vectorization 到 LSX lowering 的关键链路、正式结果、重复测量验证
以及与完整 SIMD 覆盖候选版的边界比较。

\begin{figure}[htbp]
    \centering
    \includegraphics[width=0.97\textwidth]{intro/research-mainline-roadmap.pdf}
    \caption{本文从平台缺口到主结论的研究主线与章节推进关系}
    \label{fig:research-mainline-roadmap}
\end{figure}

图~\ref{fig:research-mainline-roadmap} 将本文的主线进一步压缩为一条清晰路径：先回答 LoongArch 平台
ART 执行能力能否建立，再回答如何以统一实验设计证明其正确性与可比性，随后从候选方向中筛选出值得保留的
局部案例，最后把全文最强结论集中到 LoongArch64 最小可用 LSX 浮点 SIMD 路径上。后续章节安排也据此展开，
避免把验证细节、候选方向和主结论写成同一层级。

\subsection{主要贡献}

本文的主要贡献如下。

\begin{enumerate}
    \item 面向 AOSP 15 完成了 LoongArch 平台 ART 执行引擎关键模块的适配，覆盖解释器、Nterp、
    JIT 编译器后端、JNI 桥接和运行时快速路径等核心子系统，为 LoongArch 平台建立了可运行的
    ART 执行环境。
    \item 设计了面向目标设备的统一实验设计与证据组织方法，建立了覆盖目标机回归、libcore 测试和
    基准测试评测的实验框架，并以 \texttt{Snapshot-A} 至 \texttt{Snapshot-E} 统一约束正文结果、
    重复测量、正交案例与补充消融材料之间的边界。
    \item 在统一实验设计与证据组织下识别并比较了三类优化问题。其中，JIT 进入时机与运行时热路径数据搬移分别
    作为控制流与数据流方向的代表性局部案例，DEX 校验阶段链作为启动期正交案例，用于说明局部优化的
    机制、收益方向与适用边界。
    \item 补齐了 LoongArch64 平台上 ART 的最小 LSX 浮点向量 lowering 链路，并在规则浮点数值
    内核上取得最高 \texttt{+45.70\%} 的稳定性能收益；通过正式对比、重复测量验证以及与完整 SIMD
    覆盖候选版的边界比较，说明只有当优化机制与目标工作负载的主导热点相匹配时，局部改动才更可能
    转化为稳定且幅度较大的端到端收益。
\end{enumerate}

\section{论文组织结构}

全文共分为七章，各章内容安排如下。

\begin{enumerate}
    \item 第1章为绪论，介绍本文的研究背景、研究意义、国内外研究现状、研究内容、主要贡献以及
    论文整体结构。
    \item 第2章为相关技术与问题分析，主要介绍 ART 运行时体系结构、执行与编译机制、LoongArch
    架构特性以及 AOSP 15 ART 移植面临的关键技术难点。
    \item 第3章为 LoongArch 平台 ART 执行能力构建，重点说明解释器、Nterp、JIT 编译器后端、
    JNI 编译器、运行时入口点与标准库基础能力如何共同构成可运行执行环境。
    \item 第4章为实验设计与正确性验证，给出实验环境、验证层次、实验快照体系以及性能评测设计，
    为后续性能结论建立统一证据基础。
    \item 第5章为优化问题建模与代表性案例分析，说明候选方向如何完成收缩，并围绕 JIT 进入时机、
    运行时数据批量搬移和 DEX 校验阶段链讨论执行期与启动期路径中的代表性案例。
    \item 第6章为主结论与综合讨论，给出 LoongArch64 最小可用 LSX 浮点 SIMD 路径、从 IR 到目标
    汇编 lowering 的关键链路，以及三项正式方案的统一比较与叠加效应。
    \item 第7章为总结与展望，归纳全文得到的主要研究结论、系统性收益、适用边界与后续研究方向。
\end{enumerate}
%---------------------------------------------------------------------------%
